\chapter{置換積分}\label{章:置換積分}

$x(x^2+1)^{20}$ を積分しようとすると,展開によって多くの項が生じる.しかし,式の中で繰り返し使われているまとまりは $x^2+1$ だけである.それを $t$ と書けば,べきの部分は $t^{20}$ になる.残った $x$ は邪魔なのだろうか.実は $x^2+1$ の微分は $2x$ なので,この $x$ も変数を換えるために必要な形になっている.

置換積分では,このように式の一部分を新しい変数にして,合成関数の微分を逆に使う.計算だけでなく,どの区間をどの区間に対応させたのかも確かめる.特に平方根や三角関数から元の変数を表すときには,一つの値に一つの値が対応する範囲を選ぶ必要がある.

\section{置換積分の公式}\label{節:合成形の置換}

前章の一次式の積分では,内側の $ax+b$ の微分が定数 $a$ だったので,係数を直すだけで済んだ.では,内側が一般の関数 $g(x)$ のときはどうなるだろうか.ここで合成関数の微分を思い出そう.\,$F'(t)=f(t)$ であれば,
\[
 \frac{\dd}{\dd x}F(g(x))=f(g(x))g'(x)
\]
となる.この右辺の形が積分される式に現れたら,原始関数は左辺の $F(g(x))$ である.つまり,\,$f(g(x))g'(x)$ の積分は,\,$f$ の原始関数 $F$ に帰着する.これが置換積分の根拠である.

ここで,微分可能で,その導関数も連続である関数を\textbf{連続微分可能}という.\label{用語:連続微分可能}

\begin{theorem}{合成形の置換積分}
区間 $I$ で $g$ は連続微分可能とし,\,$g$ の値を含む区間で $f$ は連続とする.\,$F'=f$ ならば
\[
 \int f(g(x))g'(x)\dd x=F(g(x))+C.
\]
さらに $a,b\in I$ に対して
\[
 \int_a^b f(g(x))g'(x)\dd x
 =\int_{g(a)}^{g(b)}f(t)\dd t.
\]
この2つの式に,\,$g$ が単調であるという条件は要らない.
\end{theorem}

不定積分の式は,上で確かめた合成関数の微分そのものである.定積分の式も,基本定理によって両辺がともに
\[
 F(g(b))-F(g(a))
\]
となることから分かる.\,$g$ が途中で増加から減少へ変わっても,この証明は変わらない.後で逆関数を使う置換と比べるため,この点を覚えておこう.

実際の計算では,\,$t=g(x)$ と置き,\,$\dd t=g'(x)\dd x$ と書く.この記号は,\,$g'(x)$ と $\dd x$ を含む積を,\,$t$ による積分へ書き換えることを表している.微小な数を自由に割り算するという説明に頼る必要はなく,根拠は上の合成関数の微分にある.

\begin{bookpoint}
\textbf{POINT}\par
まとまり $g(x)$ を $t$ と置き,その微分 $g'(x)\dd x$ が残りの式にあることを確かめて $\dd t$ に換え,\,$t$ で積分して $x$ に戻す.
\end{bookpoint}

\begin{example}{合成関数の積分}
$\int x(x^2+1)^{20}\dd x$ を求めよ.
\end{example}

\paragraphbf{方針}
$t=x^2+1$ と置くと $\dd t=2x\dd x$ であり,残りの $x\dd x$ が $\frac12\dd t$ になる.

\paragraphbf{解答}
$\frac{(x^2+1)^{21}}{42}+C$

\paragraphbf{解法}
$t=x^2+1$ と置くと,\,$\dd t=2x\dd x$ だから,\,$x\dd x=\frac12\dd t$ とまとめて書き換えられる.したがって
\begin{align*}
 \int x(x^2+1)^{20}\dd x
 &=\frac12\int t^{20}\dd t\\
 &=\frac{t^{21}}{42}+C\\
 &=\frac{(x^2+1)^{21}}{42}+C
\end{align*}
となる.新しい変数 $t$ で積分した後,不定積分の答えは元の変数 $x$ に戻す.整式なので,実数全体で成り立つ.右辺を微分すると
\[
 \frac{21}{42}(x^2+1)^{20}\cdot2x
 =x(x^2+1)^{20}
\]
となり,元の式に戻る.

\paragraphbf{$x$ で割らずに置き換える}
解法では,\,$\dd x=\frac{\dd t}{2x}$ として $x$ で割ることはしなかった.\,$x\dd x$ 全体を $\frac12\dd t$ に換えたので,\,$x=0$ を除く必要がなく,答えは $x=0$ でも成り立つ.また $x^2+1$ は実数全体では単調でないが,原始関数の微分による確認には問題がない.定理に単調性の条件がないことと一致している.

\paragraphbf{まとまりの選び方}
新しい変数を選ぶときは,「見た目が複雑な部分」だけでなく,「その部分を微分したものが残りにあるか」を調べる.たとえば
\[
 \int\frac{\log x}{x}\dd x
\]
では,\,$t=\log x$ と置けば $\dd t=\frac1x\dd x$ であり,\,$\int t\dd t$ に変わる.したがって,\,$x>0$ で答えは $\frac12(\log x)^2+C$ となる.複雑に見える対数を変数にすると,同時に分母の $x$ もなくなるのだ.
%
\section{分母とその微分}\label{節:対数型の置換}

前節の公式で,特に $f(t)=\frac1t$ の場合を考えよう.一般的な形は
\[
 \int\frac{g'(x)}{g(x)}\dd x
\]
で与えられる.分子が分母の微分になっている分数である.\,$t=g(x)$ と置けば $\frac1t$ の積分になるので,対数に帰着させることを目標とする.

ここで1つ断りをいれておこう.本節の式は,\,$g(x)\ne0$ である区間で使う.\,$g$ が連続ならば,その区間内で符号を変えるには $0$ を通る必要があるから,符号は一定である.分母を $0$ にする点で区間が分かれ,任意定数は各区間で独立に選んでよい.

$\frac1t$ の原始関数は $\log|t|$ だから,\,$t=g(x)$,\,$\dd t=g'(x)\dd x$ と置いて
\begin{align*}
 \int\frac{g'(x)}{g(x)}\dd x
 &=\int\frac1t\dd t\\
 &=\log|t|+C\\
 &=\log|g(x)|+C
\end{align*}
を得る.実際,\,$g(x)>0$ の区間なら $\log g(x)$ を,\,$g(x)<0$ の区間なら $\log(-g(x))$ を微分すれば,どちらも同じ商 $\frac{g'(x)}{g(x)}$ に戻る.分子が $g'(x)$ の定数倍 $kg'(x)$ ならば,答えも $k\log|g(x)|+C$ となる.

\begin{bookpoint}
\textbf{POINT}\par
分子が分母 $g(x)$ の微分 $g'(x)$ の定数倍なら,\,$t=g(x)$ と置いて $\log|g(x)|$ の定数倍とする.
\end{bookpoint}

\begin{example}{絶対値を残すべき対数}
$\int\frac{6x}{3x^2-2}\dd x$ を求めよ.
\end{example}

\paragraphbf{方針}
分母 $3x^2-2$ の微分が分子の $6x$ なので,\,$t=3x^2-2$ と置いて対数にする.

\paragraphbf{解答}
$\log|3x^2-2|+C$

\paragraphbf{解法}
$t=3x^2-2$ と置けば $\dd t=6x\dd x$ なので,
\begin{align*}
 \int\frac{6x}{3x^2-2}\dd x
 &=\int\frac1t\dd t\\
 &=\log|t|+C\\
 &=\log|3x^2-2|+C
\end{align*}
となる.分母が $0$ になる $x=\pm\sqrt{\frac23}$ は除き,\,$x<-\sqrt{\frac23}$,\,$-\sqrt{\frac23}<x<\sqrt{\frac23}$,\,$\sqrt{\frac23}<x$ の各区間で考える.真ん中の区間では $3x^2-2<0$ なので,絶対値を外して $\log(3x^2-2)$ とすると実数の式ではなくなる.一方,\,$\log|3x^2-2|$ は3つの区間のどこでも使える形である.

\paragraphbf{部分分数との比較}
分母が二次式だから必ず部分分数にする,という決め方はしない.この例では分子を見れば分母の微分があり,1回の置換で終わる.前章の式の整理と置換積分は競合する規則ではなく,式の特徴に応じて選ぶ道具である.

ただし,\,$\int\frac1{3x^2-2}\dd x$ の分子は $6x$ の定数倍ではないから,同じ公式をそのまま使うことはできない.分母だけを $t$ と換えて,\,$\dd x$ を変えずに残すと,変数の異なる式を混ぜたことになる.置換後は,積分される式と微分の記号の両方を,新しい変数で書き直す必要がある.この積分は,分母を $(\sqrt3x-\sqrt2)(\sqrt3x+\sqrt2)$ と因数分解すれば,前章の部分分数で求められる.
%
\section{積分区間の変換}\label{節:置換と積分区間}

不定積分では最後に $x$ へ戻したが,定積分ならば新しい変数のまま計算を終えられる.根拠は\sectionref{節:合成形の置換}の定理の後半
\[
 \int_a^b f(g(x))g'(x)\dd x
 =\int_{g(a)}^{g(b)}f(t)\dd t
\]
である.\,$t=g(x)$ と置いたら,\,$x=a$ に対応する $t=g(a)$ を下端に,\,$x=b$ に対応する $t=g(b)$ を上端にする.大きい数を自動的に上端にするのではなく,元の端点との対応を保つ.そのため,計算の初めに $x$ と $t$ の端点の対応を
\[
 \begin{array}{c|ccc}
 x&a&\to&b\\ \hline
 t&g(a)&\to&g(b)
 \end{array}
\]
のような表にして書いておくとよい.

\begin{bookpoint}
\textbf{POINT}\par
$t=g(x)$ と置いたら端点の対応表を書き,\,$g(a)$ を下端,\,$g(b)$ を上端にして,\,$x$ に戻さずに $t$ のまま計算する.
\end{bookpoint}

\begin{example}{途中で折り返す置換}
$\int_{-1}^{2}2x\cos(x^2)\dd x$ を求めよ.
\end{example}

\paragraphbf{方針}
$t=x^2$ と置くと $\dd t=2x\dd x$ である.端点 $x=-1,2$ を $t=1,4$ に移して,\,$t$ のまま計算する.

\paragraphbf{解答}
$\sin4-\sin1$

\paragraphbf{解法}
$t=x^2$ と置けば $\dd t=2x\dd x$ であり,端点の対応は
\[
 \begin{array}{c|ccc}
 x&-1&\to&2\\ \hline
 t&1&\to&4
 \end{array}
\]
である.したがって合成形の置換積分から
\begin{align*}
 \int_{-1}^{2}2x\cos(x^2)\dd x
 &=\int_1^4\cos t\dd t\\
 &=\sin4-\sin1
\end{align*}
を得る.角 $1,4$ はラジアンである.\,$-1\le x\le2$ で $t=x^2$ は単調ではないが,定理に単調性の条件はないので,端点だけを見たこの計算で正しい.

\paragraphbf{折り返しの確認}
端点だけを見た式が正しいのは,途中の折り返しが定積分の向きに反映されているためである.図\ref{図:置換の折返し}に対応を示す.\,$x$ が $-1$ から $0$ まで増える間に $t$ は $1$ から $0$ へ減り,その後 $4$ へ増える.\,$x=0$ で積分を分けて確かめると,
\begin{align*}
 \int_{-1}^{2}2x\cos(x^2)\dd x
 &=\int_1^0\cos t\dd t+\int_0^4\cos t\dd t\\
 &=-\int_0^1\cos t\dd t+\int_0^4\cos t\dd t\\
 &=\int_1^4\cos t\dd t
\end{align*}
となる.往復した範囲 $0\le t\le1$ の寄与が符号付きで打ち消し合っているのだ.この計算は,曲線の下の面積を重複も含めて足す操作とは違う.

\begin{figure}[htbp]
 \centering
 \begin{tikzpicture}[x=1.6cm,y=0.8cm]
  \draw[-{Stealth}] (-1.35,0)--(2.35,0) node[right] {$x$};
  \draw[-{Stealth}] (0,-0.35)--(0,4.65) node[above] {$t$};
  \draw[thick,blue,domain=-1:2,samples=65] plot (\x,{\x*\x});
  \draw[dashed] (-1,0)--(-1,1)--(0,1);
  \draw[dashed] (2,0)--(2,4)--(0,4);
  \fill (-1,1) circle (1.5pt);
  \fill (2,4) circle (1.5pt);
  \node[below] at (-1,0) {$-1$};
  \node[below right] at (0,0) {$0$};
  \node[below] at (2,0) {$2$};
  \node[left] at (0,1) {$1$};
  \node[left] at (0,4) {$4$};
  \node[right] at (1.6,1.9) {$t=x^2$};
  \draw[-{Stealth},thick] (-0.88,0.85)--(-0.55,0.41);
  \draw[-{Stealth},thick] (1.12,1.38)--(1.43,2.17);
 \end{tikzpicture}
 \caption{$x$ は増え続けるが,\,$t=x^2$ は減少してから増加する.}
 \label{図:置換の折返し}
\end{figure}

\paragraphbf{減少する置換と2つの負号}
$g$ が減少する置換では,負号が2か所に現れる.たとえば $t=1-x$ と置くと,\,$x$ が $0$ から $1$ へ動くとき,\,$t$ は $1$ から $0$ へ動く.さらに $\dd t=-\dd x$ だから,
\begin{align*}
 \int_0^1(1-x)^2\dd x
 &=-\int_1^0 t^2\dd t\\
 &=\int_0^1t^2\dd t\\
 &=\frac13
\end{align*}
となる.1行目の負号は微分 $\dd t=-\dd x$ から,2行目で負号が消えたのは上下端を入れ換えたからであり,それぞれ別の理由で現れる.途中で片方だけを落とさないよう,対応表と $\dd t$ を先に書いておくとよい.
%
\section{逆関数による置換}\label{節:逆向きの置換}

これまでは,積分される式の中に $g(x)$ とその微分 $g'(x)$ の組を探した.では,\,$\sqrt{x}$ のように根号があって,その微分が見当たらない場合はどうすればよいだろうか.この場合は,根号などのまとまりを $t=g(x)$ と置いてから,元の変数を $x=h(t)$ と解き直し,\,$x$ と $\dd x$ をすべて $t$ で書き直す.根号を含まない $t$ の式,たとえば前章の割り算や部分分数で積分できる形に帰着させることを目標とする.

ここで1つ断りをいれておこう.\,$t=g(x)$ から $x$ を $t$ の1つの関数 $x=h(t)$ として表せるのは,\,$g$ が狭義単調な区間に限る.本節では,そのような区間を選んで置換する.その理由は例題の後で述べる.

根拠は,\sectionref{節:合成形の置換}の公式を $t$ の関数 $h$ に使うことである.\,$F'=f$ ならば
\[
 \frac{\dd}{\dd t}F(h(t))=f(h(t))h'(t)
\]
だから,
\[
 \int f(h(t))h'(t)\dd t=F(h(t))+C
\]
となる.ここで $t=g(x)$ を戻すと,右辺は $F(h(g(x)))$ である.これが求める $F(x)$ になるためには,\,$h(g(x))=x$,すなわち $h$ が $g$ の逆関数であればよい.したがって,\,$x=h(t)$,\,$\dd x=h'(t)\dd t$ と書き直して $t$ で積分し,最後に $t=g(x)$ を戻せばよいのだ.

たとえば根号 $\sqrt{ax+b}$ $(a\ne0)$ を含む場合は,\,$t=\sqrt{ax+b}$ と置く.\,$t\ge0$ であり,
\[
 x=\frac{t^2-b}a,\qquad \dd x=\frac{2t}a\dd t
\]
となるから,根号が $t$ に変わる.

\begin{bookpoint}
\textbf{POINT}\par
根号などを $t$ と置いて $x=h(t)$ と解き直し,\,$\dd x=h'(t)\dd t$ として,被積分関数と $\dd x$ をすべて $t$ で書き直す.
\end{bookpoint}

\begin{example}{根号への置換}
$x>0$ において,\,$\int\frac{\sqrt{x}}{1+\sqrt{x}}\dd x$ を求めよ.
\end{example}

\paragraphbf{方針}
$t=\sqrt{x}$ と置いて $x=t^2$ と解き直し,\,$t$ の分数式にしてから割り算をする.

\paragraphbf{解答}
$x-2\sqrt{x}+2\log(1+\sqrt{x})+C$

\paragraphbf{解法}
$t=\sqrt{x}$ と置けば,\,$t>0$ で $x=t^2$,\,$\dd x=2t\dd t$ となる.分子だけでなく分母と $\dd x$ も書き換えて
\begin{align*}
 \int\frac{\sqrt{x}}{1+\sqrt{x}}\dd x
 &=\int\frac{t}{1+t}\cdot2t\dd t\\
 &=2\int\frac{t^2}{t+1}\dd t
\end{align*}
となる.根号はなくなったが分数が残った.ここで前章の割り算を使い,\,$t^2=(t+1)(t-1)+1$ と分ける.
\begin{align*}
 2\int\frac{t^2}{t+1}\dd t
 &=2\int\left(t-1+\frac1{t+1}\right)\dd t\\
 &=t^2-2t+2\log(t+1)+C
\end{align*}
$t>0$ では $t+1>0$ なので,対数の絶対値を外せる.\,$t=\sqrt{x}$ を戻すと,答えは
\[
 x-2\sqrt{x}+2\log(1+\sqrt{x})+C
\]
である.微分して確かめると
\begin{align*}
 1-\frac1{\sqrt{x}}+\frac1{(1+\sqrt{x})\sqrt{x}}
 &=1-\frac1{1+\sqrt{x}}\\
 &=\frac{\sqrt{x}}{1+\sqrt{x}}
\end{align*}
となる.置換が積分を1回で終わらせなくても,根号を有理式へ変えるという目的は達している.

\paragraphbf{枝と単調性}
この例で $x=t^2$ と書けたのは,\,$t=\sqrt{x}>0$ と決めていたからである.逆に $t=x^2$ から $x$ を求めるなら,\,$x=\sqrt{t}$ と $x=-\sqrt{t}$ の2つがある.\,$x>0$ の区間では前者,\,$x<0$ の区間では後者を使う.このように,逆関数を1つに決めるために選ぶ範囲を\textbf{枝}と呼ぶことがある.

区間上の連続関数 $g$ が逆関数を持つことと,\,$g$ が狭義単調であることは同値である.ここで狭義単調とは,\,$x_1<x_2$ ならつねに $g(x_1)<g(x_2)$ となるか,つねに $g(x_1)>g(x_2)$ となることをいう.したがって,\,$t=g(x)$ から1つの関数 $x=h(t)$ を使って戻したいなら,そのような区間を選ばなければならない.さらに逆関数の微分を $h'(t)=\frac1{g'(h(t))}$ と書く場合には,分母が $0$ でないことも必要である.

これに対し,\sectionref{節:合成形の置換}の公式では逆関数を作っていない.\,$F(g(x))$ をそのまま微分したので,単調性を仮定しなかった.定積分の公式
\[
 \int_\alpha^\beta f(h(t))h'(t)\dd t
 =\int_{h(\alpha)}^{h(\beta)}f(x)\dd x
\]
自体も同じ合成形の公式であり,\,$h$ の単調性を必要としない.単調性が必要になるのは,逆関数を1つ選んで式を書き直す段階である.この2つを区別すれば,置換に単調性が要る場合と要らない場合を混同せずに済む.
%
\section{三角置換}\label{節:三角置換の初歩}

前節では根号そのものを $t$ と置いた.しかし $\sqrt{a^2-x^2}$ $(a>0)$ では,\,$t=\sqrt{a^2-x^2}$ と置いて $x$ を解き直しても,\,$x=\pm\sqrt{a^2-t^2}$ となって根号が残ってしまう.では,根号の中の平方の差を消すにはどうすればよいだろうか.

ここで三角関数の関係式
\[
 1-\sin^2\theta=\cos^2\theta,\qquad
 1+\tan^2\theta=\frac1{\cos^2\theta}
\]
を思い出そう.\,$a^2-x^2$ は前者と,\,$a^2+x^2$ は後者と同じ形をしている.そこで $x$ を三角関数で表し,根号や分母を $\cos\theta$ の式に変えて,前章の三角関数の積分に帰着させることを目標とする.

まず $\sqrt{a^2-x^2}$ である.\,$x=a\sin\theta$ と置くと
\begin{align*}
 \sqrt{a^2-x^2}
 &=a\sqrt{1-\sin^2\theta}\\
 &=a\sqrt{\cos^2\theta}\\
 &=a|\cos\theta|
\end{align*}
となる.最後は $a\cos\theta$ ではなく,まず絶対値付きである.\,$\sqrt{z^2}=|z|$ だからである.そこで $\theta$ の範囲として $-\frac\pi2\le\theta\le\frac\pi2$ を選ぶ.この範囲では $\cos\theta\ge0$ となり,絶対値を外せる.さらにこの範囲では,図\ref{図:正弦の枝}のように $\sin\theta$ が狭義単調増加なので,\,$-a\le x\le a$ の各 $x$ に対して $\theta$ が1つに決まる.前節で述べた枝を選んだことになる.このとき
\[
 \sqrt{a^2-x^2}=a\cos\theta,\qquad
 \dd x=a\cos\theta\dd\theta
\]
であり,\,$\sqrt{a^2-x^2}\dd x=a^2\cos^2\theta\dd\theta$ となる.あとは二倍角の公式で積分できる.
次に $a^2+x^2$ である.\,$x=a\tan\theta$ と置き,\,$-\frac\pi2<\theta<\frac\pi2$ を選ぶ.この範囲で正接は狭義単調増加し,すべての実数 $x$ に $\theta$ が1つずつ対応する.このとき
\[
 a^2+x^2=a^2(1+\tan^2\theta)=\frac{a^2}{\cos^2\theta},\qquad
 \dd x=\frac a{\cos^2\theta}\dd\theta
\]
となる.分母と微分に同じ $\frac1{\cos^2\theta}$ が現れるので,
\begin{align*}
 \frac1{a^2+x^2}\dd x
 &=\frac{\cos^2\theta}{a^2}\cdot\frac a{\cos^2\theta}\dd\theta\\
 &=\frac1a\dd\theta
\end{align*}
と,被積分関数が定数になる.

\begin{figure}[htbp]
 \centering
 \begin{tikzpicture}[x=2cm,y=1.5cm]
  \draw[-{Stealth}] (-1.95,0)--(1.95,0) node[right] {$\theta$};
  \draw[-{Stealth}] (0,-1.3)--(0,1.35) node[above] {$\sin\theta$};
  \draw[thick,blue,domain=-1.570796:1.570796,samples=70]
    plot (\x,{sin(\x r)});
  \draw[dashed] (-1.570796,0)--(-1.570796,-1)--(0,-1);
  \draw[dashed] (1.570796,0)--(1.570796,1)--(0,1);
  \node[above] at (-1.570796,0) {$-\frac\pi2$};
  \node[below] at (1.570796,0) {$\frac\pi2$};
  \node[above left] at (0,0) {$0$};
  \node[left] at (0,1) {$1$};
  \node[right] at (0,-1) {$-1$};
 \end{tikzpicture}
 \caption{正弦はこの区間で狭義単調増加し,逆向きの対応を1つに決められる.}
 \label{図:正弦の枝}
\end{figure}

\begin{bookpoint}
\textbf{POINT}\par
$\sqrt{a^2-x^2}$ には $x=a\sin\theta$ $\left(-\frac\pi2\le\theta\le\frac\pi2\right)$,\,$a^2+x^2$ には $x=a\tan\theta$ $\left(-\frac\pi2<\theta<\frac\pi2\right)$ と置き,端点を $\theta$ に移して $\theta$ で積分する.
\end{bookpoint}

\begin{example}{四分円の面積}
$a>0$ とする.\,$\int_0^a\sqrt{a^2-x^2}\dd x$ を求めよ.
\end{example}

\paragraphbf{方針}
$x=a\sin\theta$ と置き,\,$0\le\theta\le\frac\pi2$ を選んで $\cos^2\theta$ の積分に帰着させる.

\paragraphbf{解答}
$\frac{\pi a^2}{4}$

\paragraphbf{解法}
ここでは $x\ge0$ だから,\,$x=a\sin\theta$ と置き,\,$0\le\theta\le\frac\pi2$ を選ぶ.端点の対応は
\[
 \begin{array}{c|ccc}
 x&0&\to&a\\ \hline
 \theta&0&\to&\frac\pi2
 \end{array}
\]
である.この範囲では $\cos\theta\ge0$ なので
\[
 \dd x=a\cos\theta\dd\theta,\qquad
 \sqrt{a^2-x^2}=a\cos\theta
\]
となり,
\begin{align*}
 \int_0^a\sqrt{a^2-x^2}\dd x
 &=a^2\int_0^{\frac\pi2}\cos^2\theta\dd\theta\\
 &=a^2\left[\frac\theta2+\frac14\sin(2\theta)\right]_0^{\frac\pi2}\\
 &=\frac{\pi a^2}{4}
\end{align*}
を得る.曲線 $y=\sqrt{a^2-x^2}$ は円 $x^2+y^2=a^2$ の上半分であり,この積分は半径 $a$ の円の4分の1の面積だから,幾何学的な結果とも一致する.積分計算そのものは,三角置換と二倍角の公式で得られている.

\begin{example}{正接への置換}
$\int_0^1\frac1{1+x^2}\dd x$ を求めよ.
\end{example}

\paragraphbf{方針}
$x=\tan\theta$ と置き,\,$0\le\theta\le\frac\pi4$ を選んで,被積分関数を定数にする.

\paragraphbf{解答}
$\frac\pi4$

\paragraphbf{解法}
$x=\tan\theta$ と置き,\,$0\le\theta\le\frac\pi4$ を選ぶ.この区間で正接は狭義単調増加し,端点 $\theta=0,\frac\pi4$ は $x=0,1$ に対応する.\,$\dd x=\frac1{\cos^2\theta}\dd\theta$ だから,
\begin{align*}
 \int_0^1\frac1{1+x^2}\dd x
 &=\int_0^{\frac\pi4}
   \frac1{1+\tan^2\theta}\cdot
   \frac1{\cos^2\theta}\dd\theta\\
 &=\int_0^{\frac\pi4}1\dd\theta\\
 &=\frac\pi4
\end{align*}
を得る.

\paragraphbf{$\theta$ の範囲を選ぶ理由}
四分円の例題で $0\le\theta\le\pi$ を選ぶと,\,$\sin\theta$ は $0$ から $1$ へ増えた後,\,$0$ へ戻る.この区間全体を $0\le x\le a$ と1対1に対応させることはできず,後半では $\cos\theta$ も負になる.枝を指定することは,こうした重複と符号を確定するために必要なのである.

\paragraphbf{不定積分にする場合}
正接の例題を不定積分として表すなら,最後に $\tan\theta=x$ を満たす $\theta$ を $x$ の関数として書くことになる.この逆関数の名前,枝,微分は\chapterref{章:置換の背景}で詳しく扱う.ここでは定積分の端点が既知の角に対応するため,逆関数の記号を導入しなくても計算を終えられる.

\subsection{図形と置換}
定積分$\int_0^1\sqrt{1-x^2}\dd x$は,半径$1$の四分円の面積だから$\frac\pi4$である.三角置換でも確かめてみよう.\,$x=\sin\theta$,\,$0\le\theta\le\frac\pi2$とすると,
\[
 \int_0^1\sqrt{1-x^2}\dd x
 =\int_0^{\frac\pi2}\cos^2\theta\dd\theta
 =\frac12\int_0^{\frac\pi2}(1+\cos2\theta)\dd\theta=\frac\pi4.
\]
この場合は図形から値を求めるほうが短い.しかし,同じ置換は$\int_0^a\sqrt{1-x^2}\dd x$にも使え,途中までの面積を表せる.変数$x$を円上の角$\theta$で表したのである.その際には,高さの$\cos\theta$に加え,横幅の変化からも$\dd x=\cos\theta\dd\theta$が現れる.円上の高さを置き換えるだけでは,面積の計算にはならない.

%
%
\section{置換後の計算}\label{節:置換後の整理}

置換を選ぶとき,「すべてがすぐ基本公式になるか」だけを条件にすると,有用な置換を見落とす.根号をなくす,指数の中身を一つにする,分母の二次式を一次式にするなど,まず何を簡単にしたいかを決める.その後で,前章の式の整理を使えばよい.

\begin{example}{置換後の整理}
$\int\frac{x^3}{1+x^2}\dd x$ を求めよ.
\end{example}

分母の $1+x^2$ を $t$ と置くと,\,$\dd t=2x\dd x$ である.分子には $x^3$ があるが,\,$x^3\dd x=x^2(x\dd x)$ と分ければ,\,$x\dd x=\frac12\dd t$ を使える.残る $x^2$ も $t-1$ と書けるので,
\[
\begin{aligned}
 \int\frac{x^3}{1+x^2}\dd x
 &=\frac12\int\frac{t-1}{t}\dd t\\
 &=\frac12\int\left(1-\frac1t\right)\dd t
 =\frac12t-\frac12\log t+C.
\end{aligned}
\]
$t=1+x^2>0$ を戻し,定数 $\frac12$ を任意定数へまとめると
\[
 \int\frac{x^3}{1+x^2}\dd x
 =\frac{x^2}{2}-\frac12\log(1+x^2)+C.
\]
$t=x^2+1$ は全実数上で一対一ではないが,残った $x^2$ を $t-1$ と書くのに逆関数 $x=h(t)$ は使っていない.\,$x$ の符号を選ぶ必要もない.

別の道もある.最初に割り算をすれば
\[
 \frac{x^3}{1+x^2}=x-\frac{x}{1+x^2}
\]
なので,第一項をそのまま積分し,第二項だけに置換を使って同じ答えを得る.どちらの場合も,式のどこを簡単にするかが分かっていれば,計算の途中で方針を見失わない.

これまでの置換では,複雑なまとまりとその微分が結び付いていた.しかし $x\e^x$ のように,二つの因子を片方の合成関数の微分としてまとめにくい積もある.この場合には,積の微分を逆に読む方法が役立つ.次章では,積分する因子と微分する因子を分けて考える.


%
\section*{章末問題}
\begin{enumerate}[label=\textbf{\thechapter.\arabic*},leftmargin=*]
 \item\label{問:置換べき} $\int(6x-4)(3x^2-4x+7)^5\dd x$ を求めよ.
 \item\label{問:置換二重対数} $\int\frac1{x\log x}\dd x$ を求めよ.実数の範囲で,どの区間に答えを適用できるかも述べよ.
 \item\label{問:置換非単調} $\int_{-2}^1 2x\e^{x^2}\dd x$ を求めよ.置換が単調でなくても成り立つ理由を説明せよ.
 \item\label{問:置換根号分母} $\int_0^3\frac1{1+\sqrt{1+x}}\dd x$ を求めよ.
 \item\label{問:置換円弧} $a>0$ とする.\,$x=a\sin\theta$ と置いて,\,$\int_0^{\frac a2}\sqrt{a^2-x^2}\dd x$ を求めよ.\,$\theta$ の区間も指定せよ.
 \item\label{問:置換正接} $x=\tan\theta$ と置いて,\,$\int_0^{\sqrt3}\frac1{1+x^2}\dd x$ を求めよ.
 \item\label{問:置換枝の違い} $t=x^2$ と置くと $x=-1,1$ がともに $t=1$ へ移る.このことから $\int_{-1}^1 1\dd x=0$ としてはいけない理由を述べよ.一方,\,$\int_{-1}^1\frac{2x}{1+x^2}\dd x=0$ は,同じ置換で正しく求められることを示せ.
\end{enumerate}

\smallskip\noindent 解答は\pageref{解答:第6章}ページ.

%
